\section*{Capítulo \ref{ch:g7:sediments-and-fossils} --- Sedimentos e fósseis}

\begin{solution}{exo:g7:sediments-and-fossils:1}
(a) \omterm{def:g7:sediments-and-fossils:fossil}{molde}; (b) \omterm{def:g7:sediments-and-fossils:fossil}{icnofóssil}; (c) \omterm{def:g7:sediments-and-fossils:sediment}{sedimento}; (d) \omterm{def:g7:sediments-and-fossils:fossil}{contramolde}; (e) \omterm{def:g7:sediments-and-fossils:layer}{estrato};
(f) \omterm{def:g7:sediments-and-fossils:fossil}{fóssil}.
\end{solution}

\begin{solution}{exo:g7:sediments-and-fossils:2}
(b) o animal morre; (d) a \omterm{def:g3:stones-pebbles-sand:sand}{lama} cobre a concha depressa; (c) novas camadas
se empilham e a concha vira pedra; (a) a \omterm{def:g5:water-shapes-the-land:erosion}{erosão} desgasta as camadas de
cima.
\end{solution}

\begin{solution}{exo:g7:sediments-and-fossils:3}
(a) Um \omterm{def:g2:water-on-earth:ocean}{mar} quente, raso e claro. (b) Águas rasas, na beira de um \omterm{def:g1:landscapes-around-us:water}{lago}, de
um \omterm{def:g1:landscapes-around-us:water}{rio} ou de um \omterm{def:g2:water-on-earth:ocean}{mar}. (c) Uma floresta de pântano ou um brejo. (d) Terra
firme, com plantas crescendo (e libélulas voando), perto de uma água onde
a \omterm{def:g3:stones-pebbles-sand:sand}{lama} se depositou.
\end{solution}

\begin{solution}{exo:g7:sediments-and-fossils:4}
(a) Falso: a maioria dos seres vivos não deixa \omterm{def:g7:sediments-and-fossils:fossil}{fóssil}; eles apodrecem ou
se dissolvem antes de serem enterrados. (b) Verdadeiro: é um \omterm{def:g7:sediments-and-fossils:fossil}{icnofóssil}.
(c) Falso: os grãos maiores e mais pesados se depositam primeiro, a \omterm{def:g3:stones-pebbles-sand:sand}{lama}
fina por último. (d) Verdadeiro.
\end{solution}

\begin{solution}{exo:g7:sediments-and-fossils:5}
(a) Arenito \qty{120}{m}, calcário \qty{180}{m}, folhelho \qty{100}{m}.
(b) $120 + 180 + 100 = \qty{400}{m}$. (c) $180 \div 400 = 0.45$:
\textbf{45~\%}. (d) O seu topo foi desgastado pela \omterm{def:g5:water-shapes-the-land:erosion}{erosão}: o que sobra
não é a camada inteira, e quanto sobra muda de um lugar para outro.
\end{solution}

\begin{solution}{exo:g7:sediments-and-fossils:6}
(a) De baixo para cima: \omterm{def:g3:stones-pebbles-sand:pebble}{cascalho}, depois \omterm{def:g3:stones-pebbles-sand:sand}{areia}, depois \omterm{def:g3:stones-pebbles-sand:sand}{lama} por cima,
embaixo da água clara. (b) Os grãos minúsculos de \omterm{def:g3:stones-pebbles-sand:sand}{lama} afundam muito
devagar: eles ainda estão caindo depois de uma hora. (c) Uma segunda
série de camadas por cima da primeira: \omterm{def:g3:stones-pebbles-sand:sand}{areia}, e depois \omterm{def:g3:stones-pebbles-sand:sand}{lama} de novo.
(d) O \omterm{def:g3:stones-pebbles-sand:pebble}{cascalho}, bem no fundo.
\end{solution}

\begin{solution}{exo:g7:sediments-and-fossils:7}
Mais chance: (b), uma concha dura, enterrada depressa embaixo de \omterm{def:g3:stones-pebbles-sand:sand}{lama}
fina, em água calma. Depois, (d): a concha é dura e fica enterrada, mas
uma cheia também pode movê-la e quebrá-la. Menos chance: (a), um corpo
mole sem partes duras, sobre pedras onde nenhuma \omterm{def:g3:stones-pebbles-sand:sand}{lama} se deposita para
cobri-lo. (c) não é muito melhor: o osso continua descoberto e é
decomposto e desgastado.
\end{solution}

\begin{solution}{exo:g7:sediments-and-fossils:8}
(a) Um oco com a forma da concha impressa nas suas paredes: um \omterm{def:g7:sediments-and-fossils:fossil}{molde}.
(b) Um \omterm{def:g7:sediments-and-fossils:fossil}{contramolde}, uma cópia de pedra da concha. (c) Não: a concha se
dissolveu; o \omterm{def:g7:sediments-and-fossils:fossil}{contramolde} é feito de \omterm{def:g6:rocks-and-minerals:mineral}{minerais} novos. (d) A água passa
facilmente entre os grãos grandes de \omterm{def:g3:stones-pebbles-sand:sand}{areia}, então ela consegue chegar
até a concha e dissolvê-la; na \omterm{def:g3:stones-pebbles-sand:sand}{lama} fina, ela passa com muito menos
facilidade.
\end{solution}

\begin{solution}{exo:g7:sediments-and-fossils:9}
(a) $1.6 \times 1000 = \qty{1600}{m}$. (b) $1600 \div 80 = 20$ camadas.
(c) $160 \div 1600 = 0.1$: \textbf{10~\%}.
\end{solution}

\begin{solution}{exo:g7:sediments-and-fossils:10}
(a) Não: um pé só deixa uma pegada num chão mole, \omterm{def:g3:stones-pebbles-sand:sand}{areia} molhada ou \omterm{def:g3:stones-pebbles-sand:sand}{lama};
o chão virou \omterm{def:g3:stones-pebbles-sand:rock}{rocha} depois. (b) A onda ou a cheia seguinte as apaga.
(c) Um novo \omterm{def:g7:sediments-and-fossils:sediment}{sedimento} precisou cobri-las e preenchê-las depressa, antes
de elas serem apagadas. (d) Um \omterm{def:g7:sediments-and-fossils:fossil}{icnofóssil}: ela mostra o que o animal fez,
e não uma parte do seu corpo.
\end{solution}

\begin{solution}{exo:g7:sediments-and-fossils:11}
(a) $20 + 5 = \qty{25}{mm}$. (b) $1000 \div 25 = 40$ pares: o testemunho
registra \textbf{40 cheias}. (c) A \omterm{def:g3:stones-pebbles-sand:sand}{areia} é pesada: ela se deposita logo
depois que o \omterm{def:g1:landscapes-around-us:water}{rio} entra no \omterm{def:g1:landscapes-around-us:water}{lago}, perto da \omterm{def:g5:water-shapes-the-land:source-mouth}{foz}; só um pouco dela vai longe.
(d) Longe da \omterm{def:g5:water-shapes-the-land:source-mouth}{foz}, só a \omterm{def:g3:stones-pebbles-sand:sand}{lama} fina ainda está na água: ela se deposita ali
devagar, quando a água está calma.
\end{solution}

\begin{solution}{exo:g7:sediments-and-fossils:12}
(a) As conchas são duras: elas duram o bastante para serem enterradas, e
muitos animais de concha viviam no fundo do \omterm{def:g2:water-on-earth:ocean}{mar}, onde a \omterm{def:g3:stones-pebbles-sand:sand}{lama} se deposita
sobre eles. (b) Os corpos moles apodrecem antes de poderem ser
enterrados. (c) Não: os animais moles quase nunca deixam \omterm{def:g7:sediments-and-fossils:fossil}{fósseis}, então a
lista feita a partir dos \omterm{def:g7:sediments-and-fossils:fossil}{fósseis} é incompleta; ela mostra sobretudo os
animais com conchas ou ossos.
\end{solution}

\begin{solution}{pb:g7:sediments-and-fossils:1}
\textbf{1.} Arenito, calcário, folhelho, giz (o primeiro embaixo da
coluna).

\textbf{2.} \qty{120}{m}, \qty{180}{m} e \qty{100}{m}; total de
\qty{400}{m}.

\textbf{3.} $180 \div 400 = 0.45$: \textbf{45~\%}.

\textbf{4.} O seu topo foi desgastado: a parte de cima da camada
desapareceu, e ninguém consegue medir o que não está mais lá.

\textbf{5.} Em águas rasas, na beira de um \omterm{def:g1:landscapes-around-us:water}{lago}, de um \omterm{def:g1:landscapes-around-us:water}{rio} ou de um \omterm{def:g2:water-on-earth:ocean}{mar}:
as marcas de ondulação se formam na \omterm{def:g3:stones-pebbles-sand:sand}{areia} em águas rasas, e as pegadas
ficam na \omterm{def:g3:stones-pebbles-sand:sand}{areia} ou na \omterm{def:g3:stones-pebbles-sand:sand}{lama} mole e molhada, muitas vezes numa praia.

\textbf{6.} Num \omterm{def:g2:water-on-earth:ocean}{mar} quente, raso e claro: os corais.

\textbf{7.} Calma: a \omterm{def:g3:stones-pebbles-sand:sand}{lama} fina só se deposita onde a água está quase parada.

\textbf{8.} Primeiro, o lugar era uma praia, com águas rasas e \omterm{def:g3:stones-pebbles-sand:sand}{areia}
mole onde animais andavam. Depois, o \omterm{def:g2:water-on-earth:ocean}{mar} o cobriu: um \omterm{def:g2:water-on-earth:ocean}{mar} quente e raso
onde cresciam corais e se empilhavam conchas. Mais tarde, a água ficou
mais calma, talvez mais funda, e só a \omterm{def:g3:stones-pebbles-sand:sand}{lama} fina se depositou, com as
conchas de amonites, animais parentes das lulas e dos polvos.

\textbf{9.} $180 \times 100 = 18\,000$~cm.

\textbf{10.} $120~\text{m} = 12\,000$~cm, e $12\,000 \times 20 =
240\,000$~anos.

\textbf{11.} $1~\text{m} = 100$~cm, e $100 \times 100 = 10\,000$~anos.

\textbf{12.} $180 \times 10\,000 = 1\,800\,000$: cerca de
\textbf{1\,800\,000 anos}, quase dois milhões de anos. Isso dá
$1\,800\,000 \div 240\,000 = 7.5$ vezes mais tempo que o arenito.
\end{solution}
